\documentclass[11pt,a4paper]{article}
\usepackage[utf8]{inputenc}
\usepackage{amsmath,amssymb,amsfonts}
\usepackage{geometry}
\usepackage{hyperref}
\usepackage{xcolor}

\geometry{margin=1in}
\definecolor{primary}{RGB}{139, 92, 246}
\definecolor{secondary}{RGB}{6, 182, 212}

\title{\textbf{\Huge Relativistic Time Dilation \& Coordinate Phase Integration}\\
\Large Theoretical Mathematics of the Time Dilation DAW Engine}
\author{\textbf{Time Dilation DAW (v4.0) Engineering Group}}
\date{\today}

\begin{document}

\maketitle

\begin{abstract}
This document details the mathematical framework governing relativistic coordinate time dilation, Lorentz velocity transforms, gravitational redshift event horizons, and digital audio phase integration in \textbf{Time Dilation DAW (v4.0)}.
\end{abstract}

\section{Introduction to Relativistic Coordinate Time ($\gamma$)}
In standard Digital Audio Workstations (DAWs), digital time progresses linearly at a constant sample rate $f_s$:
\begin{equation}
t(n) = \frac{n}{f_s}
\end{equation}
In \textbf{Time Dilation DAW}, coordinate time is modulated by a dynamic relativistic factor $\gamma(t) \in [-16.0, 16.0]$, representing localized gravitational or velocity time dilation.

\section{Velocity Time Dilation (Special Relativity)}
Under Einstein's Special Theory of Relativity, coordinate time dilation for a system moving at relative velocity $v$ with respect to speed of light $c$ is defined by the Lorentz factor:
\begin{equation}
\gamma(v) = \frac{1}{\sqrt{1 - \frac{v^2}{c^2}}}
\end{equation}
For $v \to c$, $\gamma(v) \to \infty$, causing asymptotic time dilation. In our engine, $\gamma$ is clamped to $[-16.0, 16.0]$ for real-time digital stability.

\section{Gravitational Time Dilation \& Event Horizon Redshift}
Near a massive gravitational body of Schwarzschild radius $r_s = \frac{2GM}{c^2}$, time dilation at radius $r$ follows:
\begin{equation}
\gamma(r) = \frac{1}{\sqrt{1 - \frac{r_s}{r}}} = \frac{1}{\text{redshift}}
\end{equation}
As $r \to r_s$ (approaching the Event Horizon), $\text{redshift} \to 0$, causing time stasis ($\gamma \to 0.0$).

\section{Coordinate Time Accumulation \& Phase Integration}
The cumulative relativistic coordinate time $\tau(t)$ is obtained by integrating $\gamma(t')$ over physical timeline $t$:
\begin{equation}
\tau(t) = \int_0^t \gamma(t') \, dt'
\end{equation}
In discrete time at sample rate $f_s$, the per-sample phase increment $\Delta \phi(n)$ for an oscillator at frequency $f$ is:
\begin{equation}
\Delta \phi(n) = \frac{2\pi \cdot f \cdot |\gamma(n)|}{f_s}
\end{equation}
When $\gamma(n) < 0$, the phase steps backward ($\Delta \phi < 0$), executing authentic \textit{Retrograde Time Reversal}.

\end{document}
